
\chapter{The Intake Membrane: PM4PY-LSP}
\section{A Disguised Constraint Surface}
\texttt{pm4py-lsp} is not a language server. It is a process-intelligence intake membrane disguised as a language server.

We define a seven-step pattern of Manufactured Intelligence:
\begin{enumerate}
    \item Human writes accessible language (e.g., \texttt{pm4py.discover\_dfg(df)}).
    \item The LSP observes the live language surface ($\lambda \in \mathcal{L}$).
    \item Typestate inference classifies admissibility ($O^{\star} = \alpha(O)$).
    \item Unsafe process claims ($\bot_{\mathrm{refused}}$) are diagnosed before execution.
    \item Valid intent is routed to the deterministic WASM substrate.
    \item Execution boundaries emit cryptographic receipts ($\rho$).
    \item The result becomes replayable process evidence.
\end{enumerate}

\section{LSP for Workflow-Language Constraint}
This is the core paradigm shift. When the data scientist writes code, the system treats it as a \textit{language-bearing process claim}. The membrane evaluates this claim, catches unsafe typestates, issues a diagnostic, forces a constrained repair, and finally routes the intent into the WASM-backed parity path.
